\subsection{开映射特有的性质}

命题 5. 设 X、Y 是两个拓扑空间，f 是 X 到 Y 中的一个映射，B 是 X 的拓扑的一个基。则下列命题等价：

\begin{enumerate}
\def\labelenumi{\alph{enumi})}
\item
  \(f\) 是开映射。
\item
  对每个 \(\mathbf{U} \in \mathfrak{B}\)，\(f(\mathbf{U})\) 在 \(\mathbf{Y}\) 中是开集。
\item
  对每个 \(x \in \mathbf{X}\) 及每个邻域 \(\mathbf{V}\)（\(x\) 的，\(\mathbf{X}\) 中的），\(f(\mathbf{V})\) 都是 \(f(x)\) 在 \(\mathbf{Y}\) 中的一个邻域。
\end{enumerate}

a) 与 b) 的等价性由定义及 \((\mathrm{O}_{\mathrm{I}})\) 立即得出；a) 与 c) 的等价性是 § 1 第 2 目命题 1 的推论。

命题 6. 设 R 是拓扑空间 X 上的一个等价关系；则下列三个条件等价：

\begin{enumerate}
\def\labelenumi{\alph{enumi})}
\item
  关系 \(\mathbf{R}\) 是开的。
\item
  每个关于 R 饱和的子集，其内部关于 R 是饱和的。
\item
  每个关于 R 饱和的子集，其闭包关于 R 是饱和的。
\end{enumerate}

取余集（§ 1，第 6 目，公式 (2)）可见 b) 与 c) 等价。我们来证 b) 蕴含 a)：设条件 b) 成立，并设 U 是 X 的开子集，V 是它关于 R 的饱和化；则 \(\dot{V} \supset U\)，又因按假设 \(\dot{V}\) 是饱和的，故 \(\dot{V} = V\)，从而 U 的饱和化是开集。反之，设条件 a) 成立，并设 A 是一个饱和集；若 B 是 \(\mathring{A}\) 的饱和化，则 \(\mathring{A} \subset B \subset A\)，又因按假设 B 是开集，故 \(B = \mathring{A}\)。

命题 7. 设 R 是拓扑空间 X 上的开等价关系，并设 \(\varphi: \mathbf{X} \to \mathbf{X}/\mathbf{R}\) 是典范映射。若 A 是 X 的任一关于 R 饱和的子集，则 \(\varphi(\mathbf{A})\) 在 X/R 中的闭包（相应地，内部）是 \(\varphi(\overline{\mathbf{A}})\)（相应地，\(\varphi(\mathring{\mathbf{A}})\)）。

该命题的两个断言中的每一个都可由另一个通过取余集并利用 § 1 第 6 目公式 (2) 及下述事实导出：若 B 是 X 的饱和子集，则 \(\varphi(\complement B) = \complement \varphi(B)\)。凭借命题 6，\(\overline{A}\) 是饱和的；从而 \(\varphi(\overline{A})\) 在 \(X/R\) 中是闭集，又因 \(A \subset \overline{A}\)，我们有 \(\overline{\varphi(A)} \subset \varphi(\overline{A})\)，从而 \(\varphi(A) \subset \varphi(\overline{A})\)。但由 \(\varphi\) 连续，\(\varphi(\overline{A}) \subset \overline{\varphi(A)}\)（\(\S 2\)，第 1 目，定理 1），由此即得结论。

命题 8. 设 \((\mathbf{X}_t)_{t \in \mathbf{I}}\)、\((\mathbf{Y}_t)_{t \in \mathbf{I}}\) 是以同一集合 I 为指标集的两个拓扑空间族。对每个 \(t \in \mathbf{I}\)，设 \(f_t\) 是 \(\mathbf{X}_t\) 到 \(\mathbf{Y}_t\) 中的一个开映射，并设除去有限多个指标外 \(f_t\) 都是满射。则积映射 \(f: (x_t) \to (f_t(x_t))\)（\(\prod_{t \in \mathbf{I}} \mathbf{X}_t\) 到 \(\prod_{t \in \mathbf{I}} \mathbf{Y}_t\) 中的）是开映射。

凭借命题 5，我们只须证明 \(f\) 的任意初等集 \(\prod_{t \in I} A_t\)（\(\prod_{t \in I} X_t\) 中的）的像在 \(\prod_{t \in I} Y_t\) 中是开集。而这个像是 \(\prod_{t \in I} f_t(A_t)\)，且诸假设蕴含：\(f_t(A_t)\) 在 \(Y_t\) 中是开集（对每个 \(t \in I\)），且除去有限多个指标外 \(f_t(A_t) = Y_t\)；由此即得结论。

推论. 设 \((\mathbf{X}_t)_{t \in I}\) 是拓扑空间的一个族，且对每个 \(t \in I\)，设 \(\mathbf{R}_t\) 是 \(\mathbf{X}_t\) 上的一个等价关系，并设 \(f_t\) 是典范映射 \(\mathbf{X}_t \to \mathbf{X}_t / \mathbf{R}_t\)。设 \(\mathbf{R}\) 是 \(\mathbf{X} = \prod_{t \in I} \mathbf{X}_t\) 中的等价关系“对每个 \(t \in I\)，\(\operatorname{pr}_t(x) = \operatorname{pr}_t(y) (\bmod \mathbf{R}_t)"\)

（\(x\) 与 \(y\) 之间），并设 \(f\) 是积映射 \((x_{t}) \to (f_{t}(x_{t}))\)（\(X\) 到 \(\prod_{t \in I} (X_{t}/R_{t})\) 中的）。若诸关系 \(R_{t}\) 都是开的，则关系 \(R\) 是开的，且与 \(f\) 相伴的双射是 \(X/R\) 到 \(\prod_{t \in I} (X_{t}/R_{t})\) 上的一个同胚。

R 是关系 \(f(x)=f(y)\)。由于由上面的命题 8 及 § 4 第 1 目命题 1 的推论 1，f 连续且开，故结论由第 2 目命题 3 得出。

特别地，若 R（相应地，S）是拓扑空间 X（相应地，Y）上的开等价关系，则

\[
(\mathbf {X} \times \mathbf {Y}) / (\mathbf {R} \times \mathbf {S}) \quad \text { onto } \quad (\mathbf {X} / \mathbf {R}) \times (\mathbf {Y} / \mathbf {S})
\]

\pandocbounded{\includegraphics[keepaspectratio]{images/part001_bebaa45c.jpg}}

的典范双射是一个同胚。若不假定 R 与 S 都是开的，则这个双射连续但未必是同胚，即使关系 R、S 之一是相等关系（习题 6）亦然。

